% !TeX root = ../main.tex
\newcommand{\ScoreA}{515.041450}
\newcommand{\ScoreB}{28.853901}
\newcommand{\ScoreIntercept}{590.135076}
\newcommand{\ScoreMaxRoundError}{2.123}
\newcommand{\ScoreMaxProbabilityError}{6.11e-16}
\newcommand{\ScoreCardBins}{43}

\begingroup
\setlength{\intextsep}{6pt}
\setlength{\textfloatsep}{9pt}
\setlength{\abovedisplayskip}{6pt}
\setlength{\belowdisplayskip}{6pt}
\setlength{\abovedisplayshortskip}{3pt}
\setlength{\belowdisplayshortskip}{5pt}

\subsection{刻度推导与分数分解}
设坏样本概率为 $p\in(0,1)$，好坏几率为 $O=\frac{1-p}{p}$，评分函数为 $S=A+B\ln O$。
几率翻倍时分差为 $B[\ln(2O)-\ln O]=20$，且 $A+B\ln 19=600$，得到
\begin{equation}
  B\ln 2=20,\qquad B=\frac{20}{\ln 2}\approx\ScoreB,\qquad
  A=600-B\ln 19\approx\ScoreA.
  \label{eq:score-scaling}
\end{equation}
概率与分数之间的转换为
\begin{equation}
  S(p)=A+B\ln\frac{1-p}{p},\qquad
  p(S)=\frac{1}{1+\exp\!\left(\frac{S-A}{B}\right)}.
  \label{eq:score-probability}
\end{equation}
令 $\widehat\eta_i$ 为式~\eqref{eq:woe-logistic} 的拟合对数几率。代入后可分解客户 $i$ 的总分
\begin{equation}
  S_i=A-B\widehat\eta_i=S_0+\sum_{j=1}^{8}s_j(x_{ij}),\qquad
  S_0=A-B\widehat\beta_0,\qquad s_j(x_{ij})=-B\widehat\beta_j w_j(x_{ij}).
  \label{eq:score-decomposition}
\end{equation}
评分锚点由 $O=19$ 给定，即 $p=\tfrac{1}{20}=5\%$ 时 $S=600$。令八项 WOE 均为零，
则截距取值 $b$ 对应基准概率 $p_0(b)=(1+\exp(-b))^{-1}$ 和基准分 $S_0(b)=A-Bb$，因此
\[
  \begin{aligned}
    b=0 &\quad\Longrightarrow\quad p_0(0)=50\%,\qquad S_0(0)=A\approx\ScoreA,\\
    \widehat\beta_0\approx-2.602547 &\quad\Longrightarrow\quad
    p_0(\widehat\beta_0)\approx6.8975\%,\qquad S_0(\widehat\beta_0)\approx\ScoreIntercept.
  \end{aligned}
\]
由于 $\widehat\eta_i$ 随坏样本概率严格递增，$B>0$ 时坏样本概率越高，分数越低。
对同一模型采用新刻度 $S_i'=A'-B'\widehat\eta_i$（$B'>0$），任意两位客户 $i,k$ 满足
\[
  S_i'-S_k'=-B'(\widehat\eta_i-\widehat\eta_k)=\frac{B'}{B}(S_i-S_k).
\]
截距 $A'$ 在分差中抵消，$B'/B$ 为正，分差的正、负和零均保持，因此未取整分数的先后与并列关系均保持。

\subsection{完整评分卡}
表~\ref{tab:complete-scorecard} 包含全部 \ScoreCardBins{} 档及独立基准分。
整数展示对基准分和八项档分各取最近整数，恰好半分时远离 0，再相加；
它不一定等于总分最后一次取整。计算及阈值评价均使用未取整分数，不截断分值。
表中分值未取整，仅显示三位小数，完整精度另存 \texttt{scorecard.csv}。

\begin{table}[!htbp]
  \centering
  \caption{完整评分卡}
  \label{tab:complete-scorecard}
  \scriptsize\linespread{1}\selectfont
  \setlength{\tabcolsep}{3pt}
  \renewcommand{\arraystretch}{0.96}
  \begin{minipage}[t]{0.49\linewidth}
\vspace{0pt}
\centering
\begin{tabular}{lrrrr}
\toprule
档位 & WOE & 系数 & 分值 & 整数 \\
\midrule
\multicolumn{5}{l}{\textbf{截距基准分}}\\
基准分 & --- & -2.60255 & 590.135 & 590 \\
\addlinespace[1pt]
\multicolumn{5}{l}{\textbf{额度使用率}}\\
$0$ & 0.73329 & -0.61583 & 13.030 & 13 \\
$(0,0.1]$ & 1.52938 & -0.61583 & 27.176 & 27 \\
$(0.1,0.3]$ & 0.80012 & -0.61583 & 14.217 & 14 \\
$(0.3,0.6]$ & -0.00187 & -0.61583 & -0.033 & 0 \\
$(0.6,1]$ & -1.01802 & -0.61583 & -18.089 & -18 \\
$(1,+\infty)$ & -2.08854 & -0.61583 & -37.112 & -37 \\
\addlinespace[1pt]
\multicolumn{5}{l}{\textbf{年龄}}\\
$(0,30)$ & -0.61527 & -0.40832 & -7.249 & -7 \\
$[30,40)$ & -0.44917 & -0.40832 & -5.292 & -5 \\
$[40,50)$ & -0.25234 & -0.40832 & -2.973 & -3 \\
$[50,60)$ & 0.02838 & -0.40832 & 0.334 & 0 \\
$[60,70)$ & 0.69053 & -0.40832 & 8.136 & 8 \\
$[70,+\infty)$ & 1.09995 & -0.40832 & 12.959 & 13 \\
\addlinespace[1pt]
\multicolumn{5}{l}{\textbf{30–59 天逾期次数}}\\
$0$ & 0.55345 & -0.51178 & 8.173 & 8 \\
$1$ & -0.89868 & -0.51178 & -13.271 & -13 \\
$2$ & -1.62457 & -0.51178 & -23.990 & -24 \\
$3$ & -2.04305 & -0.51178 & -30.170 & -30 \\
$\geq 4$ & -2.42905 & -0.51178 & -35.870 & -36 \\
96/98（特殊） & -2.89647 & -0.51178 & -42.772 & -43 \\
\addlinespace[1pt]
\multicolumn{5}{l}{\textbf{房地产贷款数}}\\
$0$ & -0.22430 & -0.67415 & -4.363 & -4 \\
$1$--$2$ & 0.22561 & -0.67415 & 4.388 & 4 \\
$3$ & -0.04521 & -0.67415 & -0.879 & -1 \\
$\geq 4$ & -0.62112 & -0.67415 & -12.082 & -12 \\
\addlinespace[1pt]
\bottomrule
\end{tabular}
\end{minipage}%
\hfill
\begin{minipage}[t]{0.49\linewidth}
\vspace{0pt}
\centering
\begin{tabular}{lrrrr}
\toprule
档位 & WOE & 系数 & 分值 & 整数 \\
\midrule
\multicolumn{5}{l}{\textbf{负债收入比}}\\
$0$ & -0.24666 & -0.76611 & -5.452 & -5 \\
$(0,0.4]$ & 0.16456 & -0.76611 & 3.638 & 4 \\
$(0.4,0.6]$ & -0.16705 & -0.76611 & -3.693 & -4 \\
$(0.6,1]$ & -0.46997 & -0.76611 & -10.389 & -10 \\
$(1,10]$ & -0.49359 & -0.76611 & -10.911 & -11 \\
$(10,100]$ & 0.45405 & -0.76611 & 10.037 & 10 \\
$(100,+\infty)$ & 0.12062 & -0.76611 & 2.666 & 3 \\
\addlinespace[1pt]
\multicolumn{5}{l}{\textbf{90 天及以上逾期次数}}\\
$0$ & 0.39007 & -0.53753 & 6.050 & 6 \\
$1$ & -1.95995 & -0.53753 & -30.398 & -30 \\
$2$ & -2.54814 & -0.53753 & -39.521 & -40 \\
$\geq 3$ & -3.17409 & -0.53753 & -49.230 & -49 \\
96/98（特殊） & -2.89654 & -0.53753 & -44.925 & -45 \\
\addlinespace[1pt]
\multicolumn{5}{l}{\textbf{60–89 天逾期次数}}\\
$0$ & 0.28407 & -0.35875 & 2.940 & 3 \\
$1$ & -1.81197 & -0.35875 & -18.756 & -19 \\
$\geq 2$ & -2.68735 & -0.35875 & -27.817 & -28 \\
96/98（特殊） & -2.89662 & -0.35875 & -29.984 & -30 \\
\addlinespace[1pt]
\multicolumn{5}{l}{\textbf{家属人数}}\\
$0$ & 0.15601 & -0.28302 & 1.274 & 1 \\
$1$ & -0.10192 & -0.28302 & -0.832 & -1 \\
$2$ & -0.24667 & -0.28302 & -2.014 & -2 \\
$3$ & -0.27142 & -0.28302 & -2.216 & -2 \\
$\geq 4$ & -0.46295 & -0.28302 & -3.781 & -4 \\
\addlinespace[1pt]
\bottomrule
\end{tabular}
\end{minipage}%

\end{table}

\subsection{新客户评分函数与人工案例}
设 $\boldsymbol{x}\in\mathcal X$ 为新客户的十项原始输入。用 $T$ 表示第~\ref{sec:preprocessing}~节的冻结预处理和八项入模变量选取，
令 $\widetilde{x}_j=[T(\boldsymbol{x})]_j$，$w_j(\widetilde{x}_j)$ 为按冻结档位查得的 WOE。
新客户的坏样本概率和未取整总分由下式给出。
\begin{equation}
  \begin{aligned}
    \widehat\eta(\boldsymbol{x})&=\widehat\beta_0+\sum_{j=1}^{8}\widehat\beta_j w_j(\widetilde{x}_j),\\
    \widehat p(\boldsymbol{x})&=\frac{1}{1+\exp[-\widehat\eta(\boldsymbol{x})]},\\
    S(\boldsymbol{x})&=A-B\widehat\eta(\boldsymbol{x})
      =S_0+\sum_{j=1}^{8}\bigl[-B\widehat\beta_j w_j(\widetilde{x}_j)\bigr].
  \end{aligned}
  \label{eq:score-new-customer}
\end{equation}

式中系数、分箱、WOE 和刻度均取已冻结的值。函数 \texttt{score\_customer} 接受包含十项原始字段的字典，
返回 $\widehat p(\boldsymbol{x})$、$S(\boldsymbol{x})$、整数展示分及逐项明细。

构造低、中、高三组人工输入，使用同一冻结评分卡计算。风险标签表示三个案例的相对预测风险。
表~\ref{tab:scorecard-case-detail} 展示八个入模特征的原值、档位、WOE 和档分，
另两项输入与完整计算结果见仓库的 \path{examples/scenarios/}。总分另加基准分 $S_0$。

\begin{table}[H]
  \centering
  \caption{三组人工输入的逐项评分}
  \label{tab:scorecard-case-detail}
  \footnotesize\linespread{1}\selectfont
  \setlength{\tabcolsep}{2pt}
  \renewcommand{\arraystretch}{1.12}
  \newcommand{\casevalue}[2]{\shortstack[r]{$#1$\\[1pt]$#2$}}
  \begin{tabular*}{\linewidth}{@{\extracolsep{\fill}}l*{3}{r>{$}r<{$}>{$}r<{$}}@{}}
    \toprule
    & \multicolumn{3}{c}{低风险}
    & \multicolumn{3}{c}{中风险}
    & \multicolumn{3}{c}{高风险} \\
    \cmidrule(lr){2-4}\cmidrule(lr){5-7}\cmidrule(l){8-10}
    变量 & \shortstack[r]{原值\\档位} & \text{WOE} & \text{档分}
         & \shortstack[r]{原值\\档位} & \text{WOE} & \text{档分}
         & \shortstack[r]{原值\\档位} & \text{WOE} & \text{档分} \\
    \midrule
    额度使用率 & \casevalue{0.08}{(0,0.1]} & 1.52938 & 27.176 & \casevalue{0.35}{(0.3,0.6]} & -0.00187 & -0.033 & \casevalue{0.92}{(0.6,1]} & -1.01802 & -18.089 \\
年龄 & \casevalue{58}{[50,60)} & 0.02838 & 0.334 & \casevalue{35}{[30,40)} & -0.44917 & -5.292 & \casevalue{38}{[30,40)} & -0.44917 & -5.292 \\
30--59 天逾期 & \casevalue{0}{0} & 0.55345 & 8.173 & \casevalue{1}{1} & -0.89868 & -13.271 & \casevalue{2}{2} & -1.62457 & -23.990 \\
负债收入比 & \casevalue{0.25}{(0,0.4]} & 0.16456 & 3.638 & \casevalue{0.3}{(0,0.4]} & 0.16456 & 3.638 & \casevalue{0.57}{(0.4,0.6]} & -0.16705 & -3.693 \\
90 天及以上逾期 & \casevalue{0}{0} & 0.39007 & 6.050 & \casevalue{0}{0} & 0.39007 & 6.050 & \casevalue{2}{2} & -2.54814 & -39.521 \\
房地产贷款数 & \casevalue{1}{1-2} & 0.22561 & 4.388 & \casevalue{1}{1-2} & 0.22561 & 4.388 & \casevalue{0}{0} & -0.22430 & -4.363 \\
60--89 天逾期 & \casevalue{0}{0} & 0.28407 & 2.940 & \casevalue{0}{0} & 0.28407 & 2.940 & \casevalue{1}{1} & -1.81197 & -18.756 \\
家属人数 & \casevalue{0}{0} & 0.15601 & 1.274 & \casevalue{2}{2} & -0.24667 & -2.014 & \casevalue{2}{2} & -0.24667 & -2.014 \\
\bottomrule

  \end{tabular*}
  \par\smallskip
  \parbox{\linewidth}{\scriptsize 原值在上，档位在下。档分显示三位小数，计算使用完整精度。}
\end{table}

\begin{table}[H]
  \centering
  \caption{三组人工输入的总分与概率核验}
  \label{tab:scorecard-case-summary}
  \footnotesize\linespread{1}\selectfont
  \setlength{\tabcolsep}{4pt}
  \renewcommand{\arraystretch}{1.18}
  \begin{tabular*}{\linewidth}{@{\extracolsep{\fill}}l*{6}{>{$}r<{$}}@{}}
    \toprule
    & \multicolumn{3}{c}{评分（分）} & \multicolumn{3}{c}{坏样本概率（\%）} \\
    \cmidrule(lr){2-4}\cmidrule(lr){5-7}
    案例 & \text{未取整} & \text{整数} & \text{取整误差}
         & \text{模型输出} & \text{原分反算} & \text{整数分反算} \\
    \midrule
    低风险 & 644.108 & 643 & -1.108 & 1.1283 & 1.1283 & 1.1720 \\
中风险 & 586.541 & 587 & +0.459 & 7.7415 & 7.7415 & 7.6287 \\
高风险 & 474.417 & 474 & -0.417 & 80.3443 & 80.3443 & 80.5713 \\
\bottomrule

  \end{tabular*}
\end{table}

表~\ref{tab:scorecard-case-summary} 对照模型输出与分数反算概率。未取整分数可还原模型概率，
展示分取整会带来小幅偏差。原实验在全体 150,000 条记录上核验，最大反算概率误差为 $6.11\times10^{-16}$，
逐项求和与 $A-B\eta$ 的最大差为 $1.14\times10^{-13}$ 分。
测试集最大取整绝对误差为 \ScoreMaxRoundError{} 分，低于九项各取整的 $9\times0.5=4.5$ 分理论界。
\endgroup
